
\section*{5. Connections with Expos\'e IV and consequences}
\phantomsection\label{sec:sga3-VIA-5}
\addcontentsline{toc}{section}{5. Connections with Expos\'e IV and
  consequences}

\medskip
\phantomsection\label{par:sga3-VIA-5.1}%
\noindent\textbf{5.1.---}%
\footnote{\label{note:sga3-VIA-49}Editor's note: we have changed the
title of this section, which was called ``Complements'' in the
original.}
We retain the notation of \S~3 and the hypotheses of \hyperref[thm:sga3-VIA-3.2]{Theorem~3.2}.
We then have the following commutative diagram:
\begin{center}
  \phantomsection\label{diag:sga3-VIA-5.1-right-action}%
  \begin{tikzcd}[row sep=huge,column sep=huge]
    F\times_A G\times_A G
      \arrow[r,"F\times\nu"]
      \arrow[d,shift right=0.9ex,"\operatorname{pr}_2\times G"']
      \arrow[d,shift left=0.9ex,"\lambda\times G"] &
    F\times_A G
      \arrow[d,shift right=0.9ex,"\operatorname{pr}_2"']
      \arrow[d,shift left=0.9ex,"\lambda"] \\
    G\times_A G
      \arrow[r,"\nu"]
      \arrow[d,"p\times G"'] &
    G
      \arrow[d,"p"] \\
    (F\backslash G)\times_A G
      \arrow[r,dashed,"\rho"] &
    F\backslash G
  \end{tikzcd}
\end{center}
It satisfies
\[
  \operatorname{pr}_2\circ(F\times\nu)
    =
  \nu\circ(\operatorname{pr}_2\times G),
  \qquad
  \lambda\circ(F\times\nu)
    =
  \nu\circ(\lambda\times G).
\]
Moreover, since \(G\) is flat over \(A\), the left-hand vertical
sequence is exact by \hyperref[par:sga3-VIA-4.7]{4.7}. Hence \(\nu\) induces a morphism of
\(A\)-schemes
\[
  \rho:(F\backslash G)\times_A G\longrightarrow F\backslash G.
\]

The morphism \(\rho\) makes \(G\) act on the right on
\(F\backslash G\), as is immediately checked; moreover, the
canonical morphism \(G\to F\backslash G\) commutes with the right
actions of \(G\) on \(G\) and on \(F\backslash G\).%
\footnote{\label{note:sga3-VIA-50}Editor's note: we have added what
follows.}
This proves the first assertion of \hyperref[concl:sga3-VIA-3.2-ii-prime]{Theorem~3.2}(ii\(^{\prime}\)).
By \SGARefLink{prop:sga3-VIA-2.5.4}{2.5.4}, the connected components of
\[
  X=F\backslash G
\]
are therefore of finite type, irreducible, and all of the same
dimension.

To compute this dimension, we may suppose that \(A=k\) and that
\(k\) is algebraically closed. By \SGARefLink{sga3:i:scholium:2-3-3-1}{Expos\'e~I, 2.3.3.1}, the
stabilizer of the \(k\)-point \(p(e)\) is represented by the fiber
\[
  H=p^{-1}\bigl(p(e)\bigr).
\]
Since \(F\backslash G\) is the quotient of \(G\) by \(F\) in the
category of ringed spaces, the underlying space of this fiber is
\(u(F)\). As \(\ker(u)\) is finite, we obtain
\[
  \dim H=\dim u(F)=\dim F.
\]
It follows from \hyperref[concl:sga3-VIA-2.5.4-ii]{2.5.4(ii)} that
\[
  \dim X=\dim G-\dim F.
\]
This proves \hyperref[concl:sga3-VIA-3.2-ii-prime]{Theorem~3.2}(ii\(^{\prime}\)), and hence also
\hyperref[concl:sga3-VIA-3.3.2-iv]{Theorem~3.3.2(iv)}.
